\subsection{Ebenendarstellungen umformen}
\label{subsec:analytic_geometry_converting_plane_representations}

Parameter- und Normalenform beschreiben dieselbe Ebene aus zwei Blickwinkeln.
Die Umrechnung macht jeweils die für eine Aufgabe günstigere Information
sichtbar.

\begin{center}
  \CachedTikzFigure{content/resources/chapters/analytic_geometry/tikz/plane_representation_conversion.tex}
\end{center}

\subsubsection{Parameterform in Normalenform überführen}

\begin{KeyPointBox}{Vorgehen: Parameterform in Normalenform}{analytic_geometry_parameter_to_normal_method}
    Sei eine Ebene $E$ in Parameterform gegeben:
    \[
      E:\quad
      \vec x=\overrightarrow{OP}+\lambda\vec a+\mu\vec b.
    \]
    Um die Normalenform zu erhalten, gehen wir wie folgt vor:
    \begin{enumerate}[leftmargin=1.6em,itemsep=0.25em]
      \item Stützvektor $\overrightarrow{OP}$ übernehmen.
      \item Normalenvektor bestimmen: $\vec n=\vec a\times\vec b$.
      \item Normalenform aufstellen:
      \[
      E:\quad (\vec x-\overrightarrow{OP})\cdot\vec n=0.
    \]
    \end{enumerate}
\end{KeyPointBox}

\begin{ExampleBox}{Parameterform in Normalenform überführen}{analytic_geometry_parameter_to_normal_example}
    Gegeben ist die bereits bekannte Ebene
    \[
      E:\quad
      \vec x=\begin{pmatrix}6\\3\\-5\end{pmatrix}
        +\lambda\begin{pmatrix}1\\1\\-2\end{pmatrix}
        +\mu\begin{pmatrix}-1\\3\\1\end{pmatrix}.
    \]
    Ein Normalenvektor entsteht aus dem Kreuzprodukt der Richtungsvektoren:
    \begin{ClozeArea}[34]{12}
      \[
        \vec n
        =\begin{pmatrix}1\\1\\-2\end{pmatrix}
          \times\begin{pmatrix}-1\\3\\1\end{pmatrix}
        =\begin{pmatrix}
            1\cdot1-(-2)\cdot3\\
            (-2)\cdot(-1)-1\cdot1\\
            1\cdot3-1\cdot(-1)
          \end{pmatrix}
        =\begin{pmatrix}7\\1\\4\end{pmatrix}.
      \]
    \end{ClozeArea}
    Damit lautet die Normalenform
    \begin{ClozeArea}[34]{5}
      \[
        E:\quad
        \left(\vec x-\begin{pmatrix}6\\3\\-5\end{pmatrix}\right)
          \cdot\begin{pmatrix}7\\1\\4\end{pmatrix}=0.
      \]
    \end{ClozeArea}
\end{ExampleBox}

\subsubsection{Normalenform in Parameterform überführen}

\begin{KeyPointBox}{Vorgehen: Normalenform in Parameterform}{analytic_geometry_normal_to_parameter_method}
    Sei eine Ebene gegeben in Normalenform:
    \[
      E:\quad (\vec x-\overrightarrow{OP})\cdot\vec n=0.
    \]
    Um die Parameterform zu erhalten, gehen wir wie folgt vor:
    \begin{enumerate}[leftmargin=1.6em,itemsep=0.25em]
      \item Stützvektor $\overrightarrow{OP}$ aus der Normalenform übernehmen.
      \item Wähle 2 linear unabhängige Richtungsvektoren $\vec a,\vec b$ mit
        \[
          \vec n\cdot\vec a = 0\quad\text{und}\quad \vec n\cdot\vec b=0
        \]
        \textit{Tipp:} Eine Möglichkeit, zu
        $\vec n=(n_1,n_2,n_3)^{\mathsf T}$ senkrechte Richtungsvektoren zu
        konstruieren, ist für $n_2\neq0$
        \[
          \vec a=\begin{pmatrix}-n_2\\n_1\\0\end{pmatrix},
          \qquad
          \vec b=\begin{pmatrix}0\\n_3\\-n_2\end{pmatrix}.
        \]
      \item Parameterform aufstellen:
      \[
        E:\quad \vec x=\overrightarrow{OP}+\lambda\vec a+\mu\vec b,
        \qquad \lambda,\mu\in\R.
      \]
    \end{enumerate}
  \end{KeyPointBox}

\begin{ExampleBox}{Normalenform in Parameterform überführen}{analytic_geometry_normal_to_parameter_example}
    Gegeben ist die Ebene
    \[
      E:\quad
      \left(\vec x-\begin{pmatrix}2\\1\\0\end{pmatrix}\right)
        \cdot\begin{pmatrix}2\\-1\\1\end{pmatrix}=0.
    \]
    Den Stützvektor übernehmen wir direkt. Für zwei Richtungsvektoren lösen
    wir $\vec n\cdot\vec r=0$ und wählen beispielsweise
    \begin{ClozeArea}[34]{8}
      \[
        \vec a=\begin{pmatrix}1\\2\\0\end{pmatrix},
        \qquad
        \vec b=\begin{pmatrix}1\\1\\-1\end{pmatrix},
        \qquad
        \vec n\cdot\vec a=\vec n\cdot\vec b=0.
      \]
    \end{ClozeArea}
    Somit besitzt die Ebene beispielsweise die Parameterdarstellung
    \[
      E:\quad
      \vec x=\begin{pmatrix}2\\1\\0\end{pmatrix}
        +\lambda\begin{pmatrix}1\\2\\0\end{pmatrix}
        +\mu\begin{pmatrix}1\\1\\-1\end{pmatrix},
      \qquad \lambda,\mu\in\R.
    \]
\end{ExampleBox}
